% Curs 11 - Pagina 8 (13-XII-1994).
% Fragment pentru % Curs 11 - Pagina 8 (13-XII-1994).
% Fragment pentru % Curs 11 - Pagina 8 (13-XII-1994).
% Fragment pentru % Curs 11 - Pagina 8 (13-XII-1994).
% Fragment pentru \input{courses/pages/curs11-pagina8-ro}.
% Continuarea tablei de materii de la pagina 7.
% Fara preambul si fara intrerupere fortata de pagina.

\subsection*{Curs 11 - Pagina 8}

\noindent\hspace*{2em} Signatura unei permut\u{a}ri, morfismul signatur\u{a}.
% Continuarea punctului (11) de la sfarsitul paginii 7.

\begin{enumerate}
\item[(12)] Subgrupul generat de o mul\c{t}ime de elemente. Grupuri finit generate,
grupuri ciclice, teorema de structur\u{a} a grupurilor abeliene.
\item[(13)] Subgrupul derivat al unui grup.
\item[(14)] Ac\c{t}iunea unui grup pe o mul\c{t}ime. Stabilizator, orbit\u{a},
echivalen\c{t}a claselor; ac\c{t}iunea unui grup pe o mul\c{t}ime finit\u{a},
ecua\c{t}ia claselor de elemente conjugate.
% Abrevierea din manuscris ("ec. claselor") poate fi citita si
% ca "echivalenta claselor"; a doua aparitie este in contextul conjugarii.
\item[(15)] Teorema lui Cauchy.
\item[(16)] $p$-grupuri, $p$-subgrupuri Sylow, teoremele lui Sylow
(enun\c{t} + demonstra\c{t}ie m\u{a}car la una).
\item[(17)] Teoreme de izomorfism pentru grupuri (I + II).
\item[(18)] \underline{F\u{a}r\u{a}} teorema lui Burnside.
% Punctul (18) si cuvantul "Fara" sunt evidentiate/taiate partial
% in manuscris; nu il interpretam drept un subiect obligatoriu.
\end{enumerate}
.
% Continuarea tablei de materii de la pagina 7.
% Fara preambul si fara intrerupere fortata de pagina.

\subsection*{Curs 11 - Pagina 8}

\noindent\hspace*{2em} Signatura unei permut\u{a}ri, morfismul signatur\u{a}.
% Continuarea punctului (11) de la sfarsitul paginii 7.

\begin{enumerate}
\item[(12)] Subgrupul generat de o mul\c{t}ime de elemente. Grupuri finit generate,
grupuri ciclice, teorema de structur\u{a} a grupurilor abeliene.
\item[(13)] Subgrupul derivat al unui grup.
\item[(14)] Ac\c{t}iunea unui grup pe o mul\c{t}ime. Stabilizator, orbit\u{a},
echivalen\c{t}a claselor; ac\c{t}iunea unui grup pe o mul\c{t}ime finit\u{a},
ecua\c{t}ia claselor de elemente conjugate.
% Abrevierea din manuscris ("ec. claselor") poate fi citita si
% ca "echivalenta claselor"; a doua aparitie este in contextul conjugarii.
\item[(15)] Teorema lui Cauchy.
\item[(16)] $p$-grupuri, $p$-subgrupuri Sylow, teoremele lui Sylow
(enun\c{t} + demonstra\c{t}ie m\u{a}car la una).
\item[(17)] Teoreme de izomorfism pentru grupuri (I + II).
\item[(18)] \underline{F\u{a}r\u{a}} teorema lui Burnside.
% Punctul (18) si cuvantul "Fara" sunt evidentiate/taiate partial
% in manuscris; nu il interpretam drept un subiect obligatoriu.
\end{enumerate}
.
% Continuarea tablei de materii de la pagina 7.
% Fara preambul si fara intrerupere fortata de pagina.

\subsection*{Curs 11 - Pagina 8}

\noindent\hspace*{2em} Signatura unei permut\u{a}ri, morfismul signatur\u{a}.
% Continuarea punctului (11) de la sfarsitul paginii 7.

\begin{enumerate}
\item[(12)] Subgrupul generat de o mul\c{t}ime de elemente. Grupuri finit generate,
grupuri ciclice, teorema de structur\u{a} a grupurilor abeliene.
\item[(13)] Subgrupul derivat al unui grup.
\item[(14)] Ac\c{t}iunea unui grup pe o mul\c{t}ime. Stabilizator, orbit\u{a},
echivalen\c{t}a claselor; ac\c{t}iunea unui grup pe o mul\c{t}ime finit\u{a},
ecua\c{t}ia claselor de elemente conjugate.
% Abrevierea din manuscris ("ec. claselor") poate fi citita si
% ca "echivalenta claselor"; a doua aparitie este in contextul conjugarii.
\item[(15)] Teorema lui Cauchy.
\item[(16)] $p$-grupuri, $p$-subgrupuri Sylow, teoremele lui Sylow
(enun\c{t} + demonstra\c{t}ie m\u{a}car la una).
\item[(17)] Teoreme de izomorfism pentru grupuri (I + II).
\item[(18)] \underline{F\u{a}r\u{a}} teorema lui Burnside.
% Punctul (18) si cuvantul "Fara" sunt evidentiate/taiate partial
% in manuscris; nu il interpretam drept un subiect obligatoriu.
\end{enumerate}
.
% Continuarea tablei de materii de la pagina 7.
% Fara preambul si fara intrerupere fortata de pagina.

\subsection*{Curs 11 - Pagina 8}

\noindent\hspace*{2em} Signatura unei permut\u{a}ri, morfismul signatur\u{a}.
% Continuarea punctului (11) de la sfarsitul paginii 7.

\begin{enumerate}
\item[(12)] Subgrupul generat de o mul\c{t}ime de elemente. Grupuri finit generate,
grupuri ciclice, teorema de structur\u{a} a grupurilor abeliene.
\item[(13)] Subgrupul derivat al unui grup.
\item[(14)] Ac\c{t}iunea unui grup pe o mul\c{t}ime. Stabilizator, orbit\u{a},
echivalen\c{t}a claselor; ac\c{t}iunea unui grup pe o mul\c{t}ime finit\u{a},
ecua\c{t}ia claselor de elemente conjugate.
% Abrevierea din manuscris ("ec. claselor") poate fi citita si
% ca "echivalenta claselor"; a doua aparitie este in contextul conjugarii.
\item[(15)] Teorema lui Cauchy.
\item[(16)] $p$-grupuri, $p$-subgrupuri Sylow, teoremele lui Sylow
(enun\c{t} + demonstra\c{t}ie m\u{a}car la una).
\item[(17)] Teoreme de izomorfism pentru grupuri (I + II).
\item[(18)] \underline{F\u{a}r\u{a}} teorema lui Burnside.
% Punctul (18) si cuvantul "Fara" sunt evidentiate/taiate partial
% in manuscris; nu il interpretam drept un subiect obligatoriu.
\end{enumerate}
